\subsection{Matriz comparativa del escenario 1}

\begin{center}
\small
\begin{tabularx}{\textwidth}{|>{\raggedright\arraybackslash}p{2.6cm}|L|L|L|}
\hline
\rowcolor{verdeuteq}\textcolor{white}{\textbf{Criterio}} & \textcolor{white}{\textbf{Secuencial (cascada)}} & \textcolor{white}{\textbf{Iterativo e incremental}} & \textcolor{white}{\textbf{Ágil}} \\ \hline
K1. Estabilidad y cambios & Las reglas de seguridad las definió el fabricante y se mantendrán siempre igual, por eso encajan con cascada. En cambio, el proveedor de sensores se elige en el tercer mes, así que ese requisito no se sabe desde el inicio y no encaja con cascada, que necesita establecer todo al inicio. & Iterativo encaja porque, como no se sabe qué sensor se va a usar, ese requisito se puede cambiar más adelante. & Ágil podría encajar en los reportes y las pantallas, porque aún no se sabe lo que se quiere y eso se puede cambiar después. Pero lo de seguridad no cambia, así que ágil no encaja para todo el sistema. \\ \hline
K2. Partes interesadas y disponibilidad & & & \\ \hline
K3. Tipo de sistema y criticidad & & & \\ \hline
K4. Riesgo principal & & & \\ \hline
K5. Documentación y marco normativo o contractual & & & \\ \hline
\end{tabularx}
\end{center}

\subsection{Recomendación y riesgos del escenario 1}

\begin{center}
\small
\begin{tabularx}{\textwidth}{|>{\raggedright\arraybackslash}p{4.5cm}|L|}
\hline
\rowcolor{verdeuteq}\textcolor{white}{\textbf{Elemento}} & \textcolor{white}{\textbf{Respuesta}} \\ \hline
Modelo recomendado (puro o híbrido; si es híbrido, qué parte sigue cada modelo) & \\ \hline
Razón 1 (dato del escenario, criterio y consecuencia) & \\ \hline
Razón 2 & \\ \hline
Razón 3 & \\ \hline
Por qué se descartan los otros dos modelos & \\ \hline
Tensión principal del escenario (el dato que apunta a otro modelo) & \\ \hline
Riesgo 1 del modelo elegido y su mitigación & \\ \hline
Riesgo 2 del modelo elegido y su mitigación & \\ \hline
\end{tabularx}
\end{center}
